\section{Multiplier and hybrid forms}
\label{ch:multiplier-hybrid}

This chapter collects the presentations of the cubic action in which Lagrange multipliers replace, fully or in part, the Gaussian elimination of the auxiliary fields.
The multiplier form is cubic at all orders as it stands and needs no inverse of any vertex; the hybrid keeps the Hubbard--Stratonovich elimination on a complement of the kernel of the quartic form and puts a multiplier on the kernel only.
The paper keeps the multiplier form as a footnote of its all-orders section; its gauge transformations and its Batalin--Vilkovisky structure have not been worked out.
The three blocks below are the text the paper carried before it was shortened: the multiplier form from the all-orders section, the hybrid from the sextic section, and the identification of the multiplier form with part of the bar--cobar equations from the bar--cobar section.
Only here are the derivations behind them: the decoupling of the kernel at quartic order, the properties of the multiplier form that follow from its elimination, and the hybrid at order $\gs^4$ in a finite-dimensional model.

\subsection{The multiplier form}
\label{sec:mult:form}

\status{Proven}\footnote{Ledger entry \recordref{R13}, with its annotations; no derivation record of its own.}
The statement is classical and at the level of existence: the multiplier form reproduces the original action exactly once the multipliers and the auxiliary fields are eliminated, at every order.
The ghost number of the multipliers and their reality, stated in the block, come from the analysis of the bar--cobar construction and of the reality of the tower (\cref{sec:mult:barcobar} and \cref{ch:conventions}).

The following block is the subsection of the paper's all-orders section that introduced the form, in the paper's notation.

\subsubsection{Multiplier form of the action}
\label{p:sec:all:multiplier}

When the auxiliary fields cannot be taken on quotients, the construction has a variant in which nothing is inverted, the multiplier form of the action.
  It is written in terms of one field $\Sigma_n \in \Hn{n}$ per vertex, $\Sigma_1 := \Psi$, which fills all the slots of $\vertex{n}$, and one multiplier $\Theta_k \in \Hn{k}$ per such field, defined in \cref{p:app:conventions}:
  \begin{equation}
  \label{p:eq:all:multiplier-action}
  \caction[\mathrm{mult}]
    := \frac{1}{2} \, \pair{\Psi}{Q_B \Psi}
      + \sum_{n \ge 3} \frac{\gs^{n - 2}}{n!} \, \vertex{n}(\Sigma_n)
      + \sum_{k \ge 2} \pair[k]{\Theta_k}{\Sigma_k - \Sigma_{k - 1} \otimes \Psi},
\end{equation}
  which is cubic at all orders as it stands.
  The multipliers enforce the chain of constraints:
\begin{equation}
  \label{p:eq:all:constraints}
  \Sigma_k
    = \Sigma_{k - 1} \otimes \Psi,
  \qquad
  k \ge 2,
  \qquad
  \Sigma_1
    := \Psi,
\end{equation}
whose solution is $\Sigma_k = \Psi^k$ for every $k$.
The fields $\Sigma_n$ of this form are not the auxiliary fields of the construction above: they equal $\Psi^n$ exactly on the solution, whereas $\Sigma_k^* = \Psi^k + O(\gs)$, and they fill the elementary vertices rather than the reduced ones.
The multiplier equations are these constraints, and on their solution the multiplier terms vanish and the vertex terms are those of the original action, exactly and at every order, with no inverse of any vertex.
The equations of the auxiliary fields are linear in the multipliers and fix them from the top down, at any truncation of the tower or as a formal series in $\gs$.
The equation of $\Sigma_k$ involves $\Theta_k$, the multiplier $\Theta_{k + 1}$ contracted with $\Psi$ on one strand, and the vertex $\vertex{k}$ read as a $k$-string state through the pairing.
Each multiplier is thus determined by the multiplier of the next level, with real coefficients.
We have not worked out the gauge law and the Batalin--Vilkovisky structure of the multiplier form.

The multipliers pair with the auxiliary fields, so their ghost number is fixed by the saturation of the $k$-string pairing:
\begin{equation}
  \label{p:eq:all:ghost-multiplier}
  \Ngh(\Theta_k)
    = 3 k,
  \qquad
  \Ngh(\Sigma_k)
    = 2 k,
\end{equation}
one strand of ghost number three against each strand of ghost number two.
Their coefficients commute and their states have the parity $k \bmod 2$.
Their strands are read in the shifted parity of the output strands, as the pairing with a symmetric $\Sigma_k$ requires (\cref{p:app:conventions}).
With $\Theta_k^\star = \Theta_k$ the multiplier terms are real, and with the reality conditions~\eqref{p:eq:conv:reality-fields}, $\Psi^\star = - \Psi$ and $\Sigma_k^\star = (-1)^k \, \Sigma_k$, the whole action is real at all orders.

This form is the one to use when the auxiliary fields cannot be taken on quotients.
That is the case when the vertices do not annihilate the kernel of the quartic form, or of the nested operators at the higher levels, so that the conditions of the type~\eqref{p:eq:hs8:kernel-condition} fail.
It needs no kernel condition at all, since nothing is inverted.
Its price is that no auxiliary field has a quadratic term, nothing is eliminated by a Gaussian integral, and the exchange structure of the construction, the reduced vertices with it, is absent.
Between the two, the hybrid of \cref{sec:mult:hybrid} keeps the Hubbard--Stratonovich elimination on a complement of the kernel and a multiplier on the kernel directions only, and its reduced sextic vertex depends on the complement.
The datum of this form is then the complement together with the constraint sector and the reduced vertex.
A coupling of the multiplier to $(\Sigma_2, \Psi)$ at order $\gs^3$ can move the complement dependence from the vertex to the constraint, but not remove it.
This was followed through order $\gs^4$ in a finite-dimensional model, and transfers to the string vertices under the independence assumption on the exchange terms and the invertibility of the quartic form on the complement.
Beyond this order, we have not worked out the complement dependence of the later reduced vertices.


\subsubsection{Elimination and what it implies}
\label{sec:mult:elimination}

We spell out the properties of the multiplier form that follow from its elimination, since the block states them in one sentence each.
Every term of~\eqref{p:eq:all:multiplier-action} has at most three fields: the vertex term $\vertex{n}(\Sigma_n)$ is linear in the single field $\Sigma_n$, which fills all the slots of the full vertex, and the multiplier term is bilinear in $\Theta_k$ and $\Sigma_k$ or trilinear in $\Theta_k$, $\Sigma_{k - 1}$ and $\Psi$.
The equations of the multipliers are algebraic, the constraints~\eqref{p:eq:all:constraints}, and their solution $\Sigma_k = \Psi^k$ is obtained recursively from $\Sigma_1 = \Psi$.
On this solution every multiplier term vanishes and the vertex terms become $\frac{\gs^{n - 2}}{n!} \vertex{n}(\Psi^n)$, so that the original non-polynomial action is recovered in one step, at all orders.

The quadratic form of the auxiliary sector is the off-diagonal pairing between $\Theta_k$ and $\Sigma_k$.
It is non-degenerate because the one-string pairing is, and the multi-string pairing is its $k$-th tensor power.
No condition of the kind needed by the Hubbard--Stratonovich form, a quartic form invertible on a quotient and vertices vanishing on its kernel, ever arises.
The propagator of a pair $(\Sigma_k, \Theta_k)$ is ultralocal: it contains no $L_0$, no modulus and no pole, so the pair carries no new degree of freedom.

On the worldsheet, the constraint term pairs $k$ strands with $k$ strands, that is, it is a gluing by $k$ plumbings with no modulus.
All the moduli of the theory remain inside the vertices $\vertex{n}$, which enter whole.
This is the difference with the Hubbard--Stratonovich form, which builds the moduli of $\vertex{n}$ out of lower vertices joined by inverse kernels.
The multiplier form is therefore the cheap existence statement and a clean formalism, at the price of doubling the tower of fields; the Hubbard--Stratonovich form carries the geometric content.
In field-theory language, it is the familiar device of making any action cubic by Lagrange multipliers, made explicit on the worldsheet by the multi-string pairing.

The gauge invariance of the multiplier form can be written down directly, without settling its gauge algebra.
Let $\delta \Psi$ be a gauge transformation of the original action, and choose $\delta \Sigma_k$ so that the constraints $C_k := \Sigma_k - \Sigma_{k - 1} \otimes \Psi$ are invariant, $\delta C_k = 0$.
The variation of the multiplier term is then $\sum_k \pair[k]{\delta \Theta_k}{C_k}$, and the transformation of the multipliers is fixed by
\begin{equation}
  \label{eq:mult:gauge-multiplier}
  \pair[k]{\delta \Theta_k}{\,\cdot\,}
    = - \sum_{n > k} \frac{\gs^{n - 2}}{n!} \, \vertex{n}\big( \,\cdot\, \otimes \delta(\Psi^{n - k}) \big),
\end{equation}
and the variation of the multiplier form equals the variation of the original action, exactly.
The gauge parameter enters only through $\delta \Psi$, so the reality of the transformation follows from the one of $\delta \Psi$, and $\delta \Theta_k$ can be chosen real, consistently with $\Theta_k^\star = \Theta_k$, for a gauge parameter with $\Lambda^\ddagger = - \Lambda$ and a real, equivariant pairing.
This shows that the action is invariant; the closure of the algebra of these transformations and the Batalin--Vilkovisky structure of the form have not been studied.
% omitted: the comparison with the explicit gauge law of the bar--cobar identification, which the source says was not rechecked.

\subsubsection{The kernel of the quartic form at quartic order}
\label{sec:mult:kernel-decoupling}

The multiplier form is the answer to a degenerate quartic form; at quartic order a degenerate form costs nothing, and we recall why.
Let $N_2$ be in the kernel of $\sprod{2}{2}$ on $\Hn{2}$, and decompose $\Sigma_2 = \Sigma_2' + N_2$.
Since the quartic vertex is graded symmetric, $N_2$ appears neither in the kinetic term $\vertex{4}(\Sigma_2^2)$ nor in the source term $\vertex{4}(\Sigma_2, \Psi^2)$ of the quartic action: the kernel direction is absent from the action, unsourced and non-interacting.
Taking $\Sigma_2$ in the quotient $\Hn{2} / \ker \sprod{2}{2}$, equivalently on a complement of the kernel, makes $\sprod{2}{2}$ injective by construction; the equation of $\Sigma_2$ then has the unique solution $[\Psi^2]$, and the equivalence with the original action at this order is untouched, since completing the square never involves $N_2$.
A kernel shift of $\Sigma_2$ is then a symmetry of zero action, and the quotient removes it.
In a finite-dimensional model with a degenerate symmetric quartic form, a kernel direction is indeed a shift symmetry of zero action and drops out of the equation of $\Sigma_2$.

Two qualifications go with this statement.
First, the word auxiliary changes meaning: $\Sigma_2$ carries no coupled degree of freedom, its kernel component being a decoupled spectator, which is weaker than being uniquely determined by $\Psi$ on all of $\Hn{2}$.
Second, in infinite dimensions the kernel must be closed and complemented, and for the elimination through a propagator the kernel must be stable; these regularity conditions parallel what every string field propagator assumes of $b_0$ and $L_0$, and they are where a real obstruction could still hide.

The decoupling holds at quartic order only.
From the quintic order on, the quotient is consistent if and only if the quintic vertex vanishes on the kernel, condition~\eqref{p:eq:hs5:kernel-condition}.
When it fails, the hybrid of the next section, or the multiplier form itself, takes over.

\subsection{The hybrid form at order \texorpdfstring{$\gs^4$}{gs4}}
\label{sec:mult:hybrid}

\status{Proven (finite-dimensional algebraic layer) -- Verified (finite model)}\record{068-multiplier-remedy-ogs4 (R75)}
The statement is established in a finite-dimensional model with generic tensors; for the string vertices it holds under the conditions stated at the end of the block and in \cref{sec:mult:hybrid-derivation}.
It applies when the quintic vertex does not vanish on the kernel of the quartic form, and it reduces to the quotient construction when it does.

The following block was part of the paper's sextic section, after the kernel conditions; it uses the paper's notation, with $\Theta_2$ the multiplier of \cref{sec:mult:form}.

Between the quotient and the multiplier form there is a hybrid, in which the multiplier acts on the kernel directions only and the elimination on the rest.
  We choose a complement $C$ of $\ker \sprod{2}{2}$ in $\Hn{2}$, write $\pi_N$ for the projection on the kernel along $C$ and $\vertex{5 \cdot 5}^{[C]}$ for the composite vertex~\eqref{p:eq:hs5:exchange-vertex} with $\sprod{2}{2}^{-1}$ replaced by the inverse of $\sprod{2}{2}$ restricted to $C$, extended by zero on the annihilator of $C$.
  Since the values of $\sprod{3}{2}$ no longer lie in the image of $\sprod{2}{2}$ when~\eqref{p:eq:hs5:kernel-condition} fails, this extension is where the complement enters.
  We take $\Sigma_2$ on the whole of $\Hn{2}$, with the multiplier $\Theta_2$ of \cref{p:app:conventions} paired with its kernel component only:
  \begin{equation}
  \label{p:eq:hs6:dressed-vertex-C}
  \begin{aligned}
    \dvertex{6}^{[C]}
      &
      := \vertex{6} - \frac{6}{5} \, \vertex{5 \cdot 5}^{[C]},
    \\
    \caction[\le 6]^{[C]}
      &
      := \caction[\le 5]
        + \frac{\gs^4}{6!} \, \dvertex{6}^{[C]}(\Psi \otimes \Sigma_2, \Psi \otimes \Sigma_2)
    \\ &
    \quad
        - \frac{\gs^4}{6!} \, \dvertex{6}^{[C]}(\Sigma_3 - \Psi \otimes \Sigma_2, \Sigma_3 - \Psi \otimes \Sigma_2)
        + \pair[2]{\Theta_2}{\pi_N(\Sigma_2 - \Psi^2)}.
  \end{aligned}
\end{equation}
  Once $\Theta_2$, $\Sigma_3$ and $\Sigma_2$ are eliminated, this action reproduces $S_{\le 6}$ through order $\gs^4$ for every choice of complement, as an identity in $\Psi$ off-shell.
  The classical value of the multiplier starts at order $\gs^3$, and it is nonzero exactly when the quintic vertex does not vanish on the kernel.
  The vertex $\dvertex{6}^{[C]}$ depends on the complement (\cref{p:sec:all:multiplier} and \cref{sec:mult:hybrid-derivation} discuss how a coupling of the multiplier can carry this dependence instead), but the eliminated action does not, through order $\gs^4$; the remainder of order $\gs^5$ does, until the next reduction.
  We established this in the finite-dimensional models of \cref{p:app:finite-models} for the form with $\Sigma_2$ alone, in which the two sextic terms above are replaced by the single term $\dvertex{6}^{[C]}(\Sigma_2^3)$ with the coefficient of the first one, with a kernel of dimension one and two.

  It holds for the string vertices if the exchange terms are not linear combinations of the elementary vertices and the quartic form is invertible on $C$.
  When the quintic vertex vanishes on the kernel, the inverse on $C$ agrees with the inverse on the quotient on all the values of $\sprod{3}{2}$, the dependence on $C$ disappears, and the quotient form above is recovered.
  We have not worked out the gauge transformation of the hybrid, its Batalin--Vilkovisky structure or its closure at order $\gs^5$.


\subsubsection{Derivation in the finite-dimensional model}
\label{sec:mult:hybrid-derivation}

We give the computation behind the block in the finite-dimensional model where it was done.
The one-string space is $\mathcal U = \mathbb{Q}^n$ and the $2$-string space is $\mathcal S = \operatorname{Sym}^2 \mathcal U$, of dimension $D = n (n + 1) / 2$.
The quartic form is a symmetric bilinear form $V_4$ on $\mathcal S$ with a kernel $N := \ker V_4$ of dimension $r \ge 1$, the quintic vertex a tensor $V_5$ on $\mathcal S \times \mathcal S \times \mathcal U$ symmetric in its two $\mathcal S$ slots, and the sextic vertex a symmetric trilinear form $V_6$ on $\mathcal S$.
No inner product on $\mathcal S$ is used.
In the paper's notation, $V_4(A, B) = \vertex{4}(A, B)$, $V_5(A, B, \Psi) = \vertex{5}(A, B, \Psi)$, and we write $P := \Psi \otimes \Psi$, the analogue of $\Psi^2$.
Total symmetry of $V_5$ and $V_6$ on $\mathcal U^{\otimes 5}$ and $\mathcal U^{\otimes 6}$ is not imposed and not used.
The arithmetic is exact (rationals), and the classical values are truncated power series in $\gs$.

The action, without the terms in $\Psi$ alone shared by both sides, and the target are:
\begin{equation}
  \label{eq:mult:hyb:action}
  \begin{gathered}
    S_4
      = - \frac{\gs^2}{24} \, V_4(\Sigma, \Sigma)
        + \frac{\gs^2}{12} \, V_4(\Sigma, P)
        + \frac{\gs^3}{5!} \, V_5(\Sigma, \Sigma, \Psi)
        + \frac{\gs^4}{6!} \, W_6(\Sigma^3),
    \\
    S_{\mathrm{eff}}
      = \frac{\gs^2}{4!} \, V_4(P, P)
        + \frac{\gs^3}{5!} \, V_5(P, P, \Psi)
        + \frac{\gs^4}{6!} \, V_6(P, P, P),
  \end{gathered}
\end{equation}
with $W_6$ a sextic vertex to be found.
A complement $C$ of $N$ in $\mathcal S$ is encoded by a linear map $\mu \colon \mathcal S \to \mathbb{Q}^r$ with $\mu \, \Xi = 1_r$, where $\Xi$ is a $D \times r$ matrix whose columns span $N$: then $C = \ker \mu$, and $\Xi \mu$ is the projector on $N$ along $C$, the $\pi_N$ of the block.
We write $B$ for a $D \times (D - r)$ matrix whose columns span $C$, $\pi_C := 1 - \Xi \mu$, and for a covector $w$, $w_C := B^{\mathrm{T}} w$ and $w_N := \Xi^{\mathrm{T}} w$.
The quintic vertex enters through
\begin{equation}
  W_\Psi
    := V_5(\cdot, \cdot, \Psi),
  \qquad
  v
    := V_5(\cdot, P, \Psi)
    = W_\Psi P,
\end{equation}
so that $v$ is the covector of $\sprod{3}{2}(\Psi^3)$ and $v_N$ its pairing with the kernel, nonzero exactly when the quintic vertex does not vanish on the kernel along $P$.

\paragraph{Elimination of the multiplier.}
The multiplier is put on the kernel directions only:
\begin{equation}
  \label{eq:mult:hyb:multiplier}
  S^{(C)}(\Sigma, \lambda)
    = S_4(\Sigma)
      + \gs^2 \, \lambda^{\mathrm{T}} \mu (\Sigma - P),
  \qquad
  \lambda \in \mathbb{Q}^r,
\end{equation}
so that $\gs^2 \lambda$ plays the role of $\Theta_2$.
The equation of $\lambda$ fixes the kernel component of $\Sigma$:
\begin{equation}
  \mu (\Sigma - P)
    = 0
  \quad \Longleftrightarrow \quad
  \Sigma
    = P + c,
  \qquad
  c
    = B y
    \in C.
\end{equation}
With $E := \partial_\Sigma S_4$, the equation of $\Sigma$ reads $E + \gs^2 \mu^{\mathrm{T}} \lambda = 0$.
Contracting with $\Xi$ and using $\mu \Xi = 1$ isolates the multiplier, and what remains is the equation on $C$ only:
\begin{equation}
  \lambda
    = - \gs^{-2} \, \Xi^{\mathrm{T}} E,
  \qquad
  B^{\mathrm{T}} E
    = 0.
\end{equation}
The pair is therefore solved exactly by the stationary point of $S_4$ on the affine slice $P + C$, and the multiplier term vanishes on the solution:
\begin{equation}
  F^{(C)}(\Psi)
    = \operatorname{ext}_{c \in C} S_4(P + c).
\end{equation}
Since $\Xi^{\mathrm{T}} V_4 = 0$, the kernel part of $E$ starts at order $\gs^3$, $\Xi^{\mathrm{T}} E = \frac{\gs^3}{60} \, v_N + O(\gs^4)$, and the classical multiplier is
\begin{equation}
  \lambda^*
    = - \frac{\gs}{60} \, v_N + O(\gs^2).
\end{equation}
It is switched on exactly by $v_N \neq 0$, it has no piece in $\gs^{-1}$, and $\gs^2 \lambda^*$ starts at order $\gs^3$, as the block states for $\Theta_2$.

\paragraph{The on-shell value.}
On the slice, the two quartic terms combine through $V_4(P + c, P + c) - 2 V_4(P + c, P) = V_4(c, c) - V_4(P, P)$, and the quintic and sextic terms are expanded in $c$:
\begin{equation}
  \label{eq:mult:hyb:slice}
  S_4(P + c)
    = S_{\mathrm{eff}}
      + \frac{\gs^4}{6!} \big[ W_6(P^3) - V_6(P^3) \big]
      + \frac{\gs^3}{60} \, v_6(c)
      - \frac{\gs^2}{24} \, K_6(c, c)
      + \frac{\gs^4}{6!} \, W_6(c^3),
\end{equation}
with the dressed source and the dressed quadratic form
\begin{equation}
  v_6
    := v + \frac{\gs}{4} \, W_6(\cdot, P, P),
  \qquad
  K_6
    := V_4 - \frac{\gs}{5} \, W_\Psi - \frac{\gs^2}{10} \, W_6(\cdot, \cdot, P).
\end{equation}
The restriction $K_6^{CC} := B^{\mathrm{T}} K_6 B = B^{\mathrm{T}} V_4 B + O(\gs)$ is invertible for small $\gs$, since $B^{\mathrm{T}} V_4 B$ is: if $B^{\mathrm{T}} V_4 B x = 0$, then $V_4 B x$ vanishes on $C$ and on $N$, hence $B x \in N \cap C = 0$.
The stationary point is $c^* = \frac{\gs}{5} \, B (K_6^{CC})^{-1} (v_6)_C + O(\gs^4)$, the cubic term feeding back at relative order $\gs^3$.
The value of a quadratic-plus-linear function at its stationary point is half its linear term, $\frac{\gs^3}{120} \, v_6(c^*)$, so that
\begin{equation}
  F^{(C)}
    = S_{\mathrm{eff}}
      + \frac{\gs^4}{6!} \big[ W_6(P^3) - V_6(P^3) \big]
      + \frac{\gs^4}{600} \, (v_6)_C^{\mathrm{T}} (K_6^{CC})^{-1} (v_6)_C
      + \frac{\gs^4}{6!} \, W_6(c^{*3})
      + O(\gs^{10}),
\end{equation}
where the last term is of order $\gs^7$.
This is the exact rewriting of the quintic order with every inverse taken on $C$.
To order $\gs^4$:
\begin{equation}
  \label{eq:mult:hyb:mismatch}
  F^{(C)} - S_{\mathrm{eff}}
    = \frac{\gs^4}{6!} \big[ W_6(P^3) - V_6(P^3) \big]
      + \frac{\gs^4}{600} \, \braket{v, G_C v}
      + O(\gs^5),
  \qquad
  G_C
    := B \, (B^{\mathrm{T}} V_4 B)^{-1} B^{\mathrm{T}}.
\end{equation}
At orders $\gs^2$ and $\gs^3$ the mismatch vanishes for every complement.
Without the multiplier, on the linear slice $\Sigma \in C$, the saddle at order $\gs^2$ is $\pi_C P$ and a mismatch $- \frac{\gs^3}{5!} \big[ V_5(P, P, \Psi) - V_5(\pi_C P, \pi_C P, \Psi) \big]$ remains already at order $\gs^3$.

\paragraph{The generalized inverse.}
The operator $G_C$ is the inverse of the quartic form restricted to $C$, extended by zero on the annihilator $C^0$ of $C$: the operator the block uses in $\vertex{5 \cdot 5}^{[C]}$.
Since $B^{\mathrm{T}} V_4 x = B^{\mathrm{T}} V_4 \pi_C x$ for every $x$, it satisfies
\begin{equation}
  V_4 G_C V_4
    = V_4,
  \qquad
  G_C V_4 G_C
    = G_C,
  \qquad
  G_C V_4
    = \pi_C,
  \qquad
  \operatorname{im} G_C
    = C,
  \qquad
  \ker G_C
    = C^0,
  \qquad
  G_C^{\mathrm{T}}
    = G_C.
\end{equation}
The first two identities and the image alone do not fix it: the reflexive inverses with image $C$ form a family indexed by their kernel, and they give different values of $\braket{v, G v}$ exactly when $v_N \neq 0$.
The kernel $C^0$ fixes $G_C$ in this family.
In particular, the Moore--Penrose inverse of a coordinate representation is one member, for the coordinate-orthogonal complement, with no distinguished role: the string field carries no auxiliary inner product.

\paragraph{The dressed vertex.}
With the composite $\mathcal C^{(C)}$, the exchange vertex $\vertex{5 \cdot 5}$ of the paper with $V_4^{-1}$ replaced by $G_C$, we have $\mathcal C^{(C)}(P, P, P) = \braket{v, G_C v}$, and the matching condition~\eqref{eq:mult:hyb:mismatch} at order $\gs^4$ is solved by
\begin{equation}
  \label{eq:mult:hyb:W6C}
  W_6^{(C)}
    = V_6 - \frac{6}{5} \, \mathcal C^{(C)},
\end{equation}
the $\dvertex{6}^{[C]}$ of the block.
As for the quotient construction, the condition fixes $W_6$ on the diagonal $P^3$ only, and~\eqref{eq:mult:hyb:W6C} is the choice with the strand structure of the exchange vertex.
The model confirms that with $W_6^{(C)}$ the mismatch at order $\gs^4$ vanishes for every complement tried (two random complements and a one-parameter family), and that it does not vanish with a coefficient other than $\frac{6}{5}$, with $V_6$ alone, with the composite of another complement, or with the Moore--Penrose composite in a non-orthogonal complement; without the multiplier, the order $\gs^3$ fails even with $W_6^{(C)}$.
For $r = 0$, $G_C = V_4^{-1}$ and the quotient construction is recovered verbatim.

\paragraph{Complement dependence.}
Fix a reference complement $C_0$ with basis $B_0$; every complement is $C_L$, spanned by the columns of $B_0 + \Xi L$, for a unique $L \in \operatorname{Hom}(C_0, N)$.
Since $V_4 \Xi = 0$, the restricted form does not depend on $L$, and only the restricted covectors move:
\begin{equation}
  \label{eq:mult:hyb:L}
  (B_0 + \Xi L)^{\mathrm{T}} V_4 (B_0 + \Xi L)
    = B_0^{\mathrm{T}} V_4 B_0
  \quad \Longrightarrow \quad
  \braket{w, G_{C_L} w'}
    = (w_{C_0} + L^{\mathrm{T}} w_N)^{\mathrm{T}} \, (V_4^{C_0 C_0})^{-1} \, (w'_{C_0} + L^{\mathrm{T}} w'_N).
\end{equation}
Hence $\braket{v, G_{C_L} v}$ is a quadratic polynomial in $L$, constant if and only if $v_N = 0$.
For the converse, when $v_N \neq 0$, the covector $y := L^{\mathrm{T}} v_N$ sweeps all of $C_0^*$ as $L$ varies, and a non-degenerate quadratic form is not constant on an affine line $v_{C_0} + t y$ with $y$ non-null.
As an identity in $\Psi$, the condition is $V_5(N, \Psi^{\otimes 2}, \Psi) \equiv 0$, which for a totally symmetric $V_5$ is the vanishing of the quintic vertex on the kernel, condition~\eqref{p:eq:hs5:kernel-condition}.
Without total symmetry, the condition on the diagonal can hold while the full condition fails; then the diagonal is complement-independent and the off-diagonal tensor is not.

When the quintic vertex vanishes on the kernel, the whole tensor $\mathcal C^{(C)}$ is complement-independent.
Indeed, on the image of $V_4$, which is the annihilator of $N$, $\braket{V_4 x, G_C V_4 x'} = V_4(x, \pi_C x') = V_4(x, x')$, so $G_C$ restricted there is the inverse of the quartic form on the quotient $\mathcal S / N$ for every $C$, and the quotient construction is recovered.

\paragraph{The next order.}
Expanding $(K_6^{CC})^{-1}$ and $v_6$ one order further, the remainder at order $\gs^5$ with $W_6^{(C)}$ is
\begin{equation}
  [\gs^5] \big( F^{(C)} - S_{\mathrm{eff}} \big)
    = \frac{1}{600} \Big[
      \frac{1}{2} \, W_6^{(C)}(G_C v, P, P)
      + \frac{1}{5} \, \braket{v, G_C W_\Psi G_C v}
    \Big],
\end{equation}
the remainder of the quotient construction at this order with $V_4^{-1}$ replaced by $G_C$.
It depends on $C$, until the next reduction.
The reduction at order $\gs^5$ is therefore expected to proceed with the same replacement, but this has not been checked: it involves $\Sigma_3$ and its own kernel questions.

\paragraph{Moving the complement dependence.}
The vertex $W_6^{(C)}$ depends on the complement, and no complement-independent sextic vertex serves two complements with $\braket{v, (G_C - G_{C'}) v} \neq 0$, within the class $\mathsf K$ where the constraint is $\mu (\Sigma - P) = 0$ undeformed, $V_5(\Sigma, \Sigma, \Psi)$ is the only coupling at order $\gs^3$ and $W_6(\Sigma^3)$ the only one at order $\gs^4$.
This class is the hybrid as proposed, continued from order $\gs^3$ with the only new coupling that order $\gs^4$ brings.
Outside it, add the cubic coupling $\gs^3 \lambda^{\mathrm{T}} Z(\Sigma, \Psi)$, with $Z \colon \mathcal S \times \mathcal U \to N$ linear in $\Sigma$ and built from $V_5$ and an inverse on a complement.
The constraint becomes $\mu (\Sigma - P) + \gs Z(\Sigma, \Psi) = 0$, and with $\Sigma = P + c + \Xi \eta$, $\eta = - \gs Z(P, \Psi) + O(\gs^2)$:
\begin{equation}
  S_4(P + c + \Xi \eta)
    = S_4(P + c) + \eta^{\mathrm{T}} \Xi^{\mathrm{T}} E + O(\eta^2)
  \quad \Longrightarrow \quad
  F^{(C, Z)} - F^{(C)}
    = - \frac{\gs^4}{60} \, v_N^{\mathrm{T}} Z(P, \Psi) + O(\gs^5),
\end{equation}
which leaves orders $\gs^2$ and $\gs^3$ untouched.
By~\eqref{eq:mult:hyb:L} with a reference complement $C_a$ and $A_a := (B_a^{\mathrm{T}} V_4 B_a)^{-1}$, the difference $\braket{v, G_{C_L} v} - \braket{v, G_{C_a} v} = v_N^{\mathrm{T}} L A_a (2 B_a + \Xi L)^{\mathrm{T}} v$ is $v_N^{\mathrm{T}}$ against a covector linear in $v$, so the choice
\begin{equation}
  Z_L(\Sigma, \Psi)
    = \frac{1}{10} \, L A_a (2 B_a + \Xi L)^{\mathrm{T}} \, V_5(\cdot, \Sigma, \Psi)
\end{equation}
makes the fixed vertex $W_6^{(C_a)}$ match every complement $C_L$ at order $\gs^4$; the model confirms it at corank one and two, and the same complements fail without $Z$.
In general the matching condition at order $\gs^4$ reads
\begin{equation}
  W_6(P^3)
    = V_6(P^3)
      - \frac{6}{5} \, \braket{v, G_C v}
      + 12 \, v_N^{\mathrm{T}} Z(P, \Psi).
\end{equation}
The complement dependence moves between the constraint and the vertex; it is not removed, since the reference complement is still a choice and no canonical one appears.
The datum of the presentation is the triple of the constraint, the coupling $Z$ and the vertex, as the block says in words.
What is complement-independent is the eliminated action through order $\gs^4$.

One reading of the dependence: a change of complement moves the fixed kernel component $\Xi \mu P$ of $\Sigma$ together with the vertex.
At quartic order kernel shifts of $\Sigma$ are an exact symmetry (\cref{sec:mult:kernel-decoupling}), and the constraint is its gauge fixing; when the quintic vertex does not vanish on the kernel, the shift is no longer a symmetry, and the vertex, or $Z$, has to move with the fixing.
A canonical choice would need further structure, for example a complement adapted to $V_5$, which nothing in the construction singles out.

\paragraph{Scope.}
The model has $n = 3$, $D = 6$, kernels of dimension one and two and the case $r = 0$, generic $V_5$ and $V_6$, and the orders $\gs^2$ to $\gs^5$ exact; the dependence on the complement was sampled on a line and two random points, and~\eqref{eq:mult:hyb:L} covers the rest algebraically.
For the string vertices, the statement needs the quartic form invertible on the quotient with a stable kernel, the kernel closed and complemented, and, for the exclusion within $\mathsf K$ only, exchange terms that are not linear combinations of the elementary vertices; the complement must also be chosen compatible with ghost number, level and momentum, which is possible since the quartic form is homogeneous and its kernel graded.
Not covered: the gauge transformation of the hybrid, its Batalin--Vilkovisky pairing, and its closure at order $\gs^5$ and beyond.

\subsection{Relation to the bar--cobar canonical equations}
\label{sec:mult:barcobar}

\status{mixed}\record{024-hs-vs-bar-cobar (R34)}
The ghost number $3 k$ of the multipliers is proven by the saturation count of the multi-string pairing, for the pairing that is the $k$-th tensor power of the one-string pairing; that the multiplier form is the Lagrange-multiplier action of the chain part of the bar--cobar equations was checked numerically in finite-dimensional models with commuting fields; that it is not the action of their full system is shown term by term, the failing term being the split $(2, 2)$.

The following block is the paragraph of the paper's bar--cobar section that identified the multiplier form, after the canonical equations of~\cite{Chiaffrino:2026:UniversalQuadraticField} were written in the variables of the cubic action (\cref{p:sec:barcobar:canonical}).
In it, \enquote{their} refers to the construction of~\cite{Chiaffrino:2026:UniversalQuadraticField}, and the doubling by multipliers is the one mentioned at the end of \cref{p:sec:barcobar:canonical}.

In these variables, we recognize the multiplier form of the cubic action, treated in \cref{p:sec:all:multiplier}:
\begin{equation}
  \label{p:eq:barcobar:multiplier-action}
  \caction[\mathrm{mult}]
    = \frac{1}{2} \, \pair{\Psi}{Q_B \Psi}
      + \sum_{n \ge 3} \frac{\gs^{n-2}}{n!} \, \vertex{n}(\Sigma_n)
      + \sum_{k \ge 2} \pair[k]{\Theta_k}{\Sigma_k - \Sigma_{k-1} \otimes \Psi},
\end{equation}
with multipliers $\Theta_k \in \Hn{k}$ (\cref{p:sec:all:multiplier}) of ghost number $3 k$ (\cref{p:app:conventions}).
Its equations of motion map onto the canonical equations~\eqref{p:eq:barcobar:canonical-eqs} as follows.
The $\Theta_k$ equations impose the conditions $\Sigma_k = \Sigma_{k-1} \otimes \Psi$, which we call the chain: these are the splits $(k - 1, 1)$ of the factorization condition, in their part symmetric under all strands.
The $\Psi$ equation, once the multipliers are eliminated through the $\Sigma_n$ equations, is the original equation $\mathcal E_\Psi = 0$, that is, the one-output component of $B \Phi_1 = 0$.
The $\Sigma_n$ equations themselves have no counterpart in their system. They are linear in the multipliers and, for the tower truncated at a finite length, determine them from the top down.

Conversely, three groups of their equations have no Euler--Lagrange counterpart.
The first two are the splits with $2 \le i \le n - 2$, the first of which is $\Sigma_4 = \Sigma_2 \otimes \Sigma_2$, and the symmetry conditions on the right-hand sides of~\eqref{p:eq:barcobar:factorization}.
The third group is the components of $B \Phi_1 = 0$ with two or more output strands.
All of them hold identically on the diagonal $\Sigma_k = \Psi^k$, to which the chain reduces: this solution satisfies every split and reduces every component of $B \Phi_1 = 0$ to the one-output one.

Hence, for the multiplier form truncated at a finite length, the solutions of its equations of motion, modulo the values of the multipliers, are in one-to-one correspondence with their canonical solutions truncated at the same length.
The correspondence follows from the triangular form of the chain, which fixes $\Sigma_k = \Psi^k$ once $\Psi$ is given.
We have checked the term-by-term map in finite-dimensional models with commuting fields, of the kind described in \cref{p:app:finite-models}.

In short, the multiplier form is the Lagrange-multiplier action of the universal part of their system only, namely the chain.
The theory enters through the kinetic term of $\Psi$ and the vertices $\vertex{n}(\Sigma_n)$, which are linear in the auxiliary fields.
The dual tower $\Theta_k$ supplies a non-degenerate pairing where their space has none, while $\Psi$ keeps its own; it does so by enlarging the field space, as the doubling by multipliers above does, restricted to the chain. It therefore does not answer their question of an action on their field space itself.

